\chapter{Le basis trade sur contrats à terme obligataires}\label{ch:s2:the-bond-futures-basis-trade}

Acheter un Treasury, vendre le contrat à terme sur celui-ci, financer l'obligation en repo et encaisser les quelques points de base dont le taux de repo implicite du contrat à terme dépasse le
taux auquel on se finance. À son sommet, avant mars 2020, l'Office of Financial Research estimait que les positions de basis trade des hedge funds atteignaient 400 à 500 milliards de dollars. Cela
représentait plus de 60~\% de leur exposition aux Treasuries et plus de 70~\% de leurs emprunts en repo. Entre le 18 février et le 17 mars 2020, ils ont réduit leurs
positions courtes sur les contrats à terme à 2, 5 et 10 ans, de 659 à 554 milliards de dollars. La Réserve fédérale s'est alors engagée à acheter des Treasuries dans
les montants nécessaires. Sur le marché synthétique de ce chapitre, un book de basis trade à 50 fois son capital rapporte 10.3~\% par an jusqu'à un stress de financement
artificiel de dix jours. Les appels de marge l'obligent à vendre au pire niveau du basis, et il achève l'épisode qui suit le stress en retrait de 9.6~\% par rapport à son point de départ. À dix fois son capital, il gagne 2.1~\%
et ne termine qu'en baisse de 0.2~\%. La construction est \texttt{firm.basistrade}.

\section{Mécanique~: comptant, contrat à terme, repo}

Un contrat à terme sur Treasury (livre~2, chapitre~6) est livré par le vendeur, qui choisit dans un panier l'obligation la moins chère à livrer (CTD). La base brute est le
prix de l'obligation moins le prix du contrat à terme multiplié par son facteur de concordance. Nette du carry jusqu'à la livraison, elle devient la base nette. Le taux de repo implicite est le taux de financement
auquel l'achat de l'obligation, sa livraison au contrat à terme et le remboursement de l'emprunt s'équilibrent. Lorsque le taux de repo implicite dépasse le taux de repo, une
position longue sur la base gagne l'écart~: longue sur l'obligation financée en repo, courte sur le contrat à terme. Le book ci-dessous gagne 20 points de base par an sur le nominal
qu'il détient, hypothèse propre à ce chapitre.

Le basis trade comporte trois sources de risque. La base peut évoluer à l'encontre du book avant la livraison. Le prêt en repo doit être renouvelé, souvent chaque nuit. Et la
bourse des contrats à terme et le prêteur en repo détiennent chacun du collatéral~: une marge sur les contrats à terme et une décote sur l'obligation.

\begin{definition}[Effet de levier du basis trade]\label{def:s2:the-bond-futures-basis-trade:leverage}
L'\emph{effet de levier du basis trade}\index{effet de levier du basis trade} est la valeur nominale des obligations détenues en position longue (et des contrats à terme en position courte) par unité de capital du trader~; elle
est plafonnée par le collatéral dont la stratégie a besoin, si bien qu'elle vaut au plus un divisé par la somme de la décote du repo et de la marge sur contrats à terme (par unité de nominal), et elle baisse dès que
l'une des deux augmente.
\end{definition}

\begin{definition}[Risque de renouvellement du repo]\label{def:s2:the-bond-futures-basis-trade:rollover}
Le \emph{risque de renouvellement du repo}\index{risque de renouvellement du repo} est le risque qu'une position financée en repo à court terme ne puisse pas être refinancée aux mêmes conditions
lorsque le prêt arrive à échéance~: à un taux plus élevé, avec une décote plus forte, pour un montant moindre, ou pas du tout.
\end{definition}

\section{Effet de levier et acteurs}

Barth et Kahn ont retracé la stratégie à partir des données réglementaires. Les sociétés de gestion voulaient de la duration sans mobiliser leur bilan, et ont acheté des contrats à terme sur Treasury.
Cela a poussé les contrats à terme au-dessus de leur prix sans arbitrage, et les hedge funds ont répondu à cette demande en prenant une position longue sur la base. De fin 2017 à septembre 2019,
l'exposition des hedge funds aux Treasuries est passée de 1.06 billion de dollars à 2.02 billions de dollars, et leurs positions courtes sur contrats à terme se sont accrues de 352 milliards de dollars. Les auteurs décrivent les fonds comme
des entrepôts. Ils conservent des Treasuries pour les détenteurs de contrats à terme longs, les financent en repo et relient ainsi les marchés des Treasuries, des contrats à terme et du repo.

\begin{dated}{2026-09}{Fonds à effet de levier sur contrats à terme sur Treasury}\label{dat:s2:the-bond-futures-basis-trade:tff}
Le rapport Traders in Financial Futures de la CFTC classe chaque grand trader dans une catégorie~; les «~fonds à effet de levier~» incluent les hedge funds et d'autres
gestions pour compte de tiers (managed money), qui détiennent des positions de natures très variées, pas seulement des basis trades. En additionnant leurs positions nettes sur les six contrats à terme sur Treasury du CBOT (2-, 5-, 10-year, Ultra
10-year, bond et Ultra bond) à leur valeur nominale, elles étaient nettes acheteuses de 34 milliards de dollars en moyenne de juillet 2010 à 2014 et nettes vendeuses de 116 milliards de dollars à la fin de
2017. Elles étaient nettes vendeuses de 475 milliards de dollars le 18 février 2020 et de 376 milliards de dollars le 17 mars 2020, avaient atteint une position nette vendeuse de 1\,184 milliards de dollars le 12 novembre
2024, et étaient nettes vendeuses de 787 milliards de dollars le 15 septembre 2026.
\end{dated}

La série publique (la \cref{fig:s2:the-bond-futures-basis-trade:tff}) montre l'encombrement (crowding) dont dépend le basis trade, et l'ampleur des flux qui
surviendraient si le basis trade se débouclait.

\begin{omfigure}
\begin{tikzpicture}
\begin{axis}[width=11cm,height=5.4cm,xlabel={\small année},ylabel={\small position nette (milliards de dollars)},tick label style={font=\footnotesize},
  xmin=2010.5,xmax=2027,ymin=-1250,ymax=250,grid=major,grid style={gray!25},table/col sep=comma,/pgf/number format/1000 sep={}]
\addplot[omThm,thick] table[x=year,y=total]{figdata/strategies-2/11-the-bond-futures-basis-trade/tff.csv};
\addplot[black,thin,domain=2010.5:2027] {0};
\end{axis}
\end{tikzpicture}

\omcaption{Position nette hebdomadaire des fonds à effet de levier sur les six contrats à terme sur Treasury du CBOT, à valeur nominale, de juillet 2010 à septembre 2026 (négatif~: position nette vendeuse).
La catégorie comprend davantage que les basis trades. Dérivée des fichiers Traders in Financial Futures de la CFTC. Données~:
\texttt{s2\_fetch\_tff}.}\label{fig:s2:the-bond-futures-basis-trade:tff}
\end{omfigure}

\section{Mars 2020}

Barth et Kahn décrivent le stress. Lors de la première semaine de mars 2020, les fourchettes sur les Treasuries à plus longue échéance se sont élargies, les taux de repo sur collatéral Treasury ont grimpé,
et les spreads d'arbitrage ont divergé. Les ventes des investisseurs real money ont fait monter les marges sur contrats à terme et rendu le repo volatil. Les hedge funds ont partiellement débouclé leurs
positions de basis trade~: leurs positions courtes sur les contrats à terme à 2, 5 et 10 ans sont passées de 659 à 554 milliards de dollars entre le 18 février et le 17 mars. Les auteurs estiment
que de grands acteurs du basis trade ont vendu entre 91 et 105 milliards de dollars de Treasuries entre fin février et fin mars. Ils font aussi valoir que le
basis trade était peu susceptible d'être la cause principale du stress avant le 17 mars, puisque les Treasuries qu'il détenait continuaient de se négocier avec une prime. Le 15 mars, le
Federal Open Market Committee a annoncé qu'il augmenterait ses avoirs de Treasuries d'au moins 500 milliards de dollars. Le 23 mars, il a indiqué qu'il les achèterait dans les
montants nécessaires pour soutenir le bon fonctionnement des marchés.

Le stress synthétique a les mêmes ingrédients. Pendant dix jours, à partir de l'année 2.5, la base évolue de 60 cents pour 100 de nominal à l'encontre du book. La décote du repo
et la marge sur contrats à terme passent chacune de 1~\% à 3~\%, et le taux de repo augmente d'un point de pourcentage. Ensuite, la base revient en deux mois,
tandis que la décote et la marge restent élevées. Un book dont le capital ne couvre plus le collatéral doit vendre jusqu'à ne plus détenir que ce que son capital finance, en payant 10 cents
pour 100 sur ce qu'il vend, et ne peut pas reprendre de levier avant la fin du stress (le \cref{lst:s2:the-bond-futures-basis-trade:book}).

\begin{center}
\small
\begin{tabular}{@{}lcccc@{}}
\toprule
cinq ans, capital 1 & 10$\times$ & 20$\times$ & 30$\times$ & 50$\times$\\
\midrule
rendement annuel avant le stress (\%) & 2.1 & 4.1 & 6.2 & 10.3\\
creux du stress (\%) & $-6.1$ & $-10.8$ & $-13.2$ & $-17.8$\\
après le rétablissement de la base (\%) & $-0.2$ & $-1.9$ & $-4.5$ & $-9.6$\\
jours avec ventes forcées & 0 & 10 & 10 & 66\\
nominal détenu après le stress & 10.4 & 16.5 & 16.8 & 17.5\\
rendement annuel sur cinq ans (\%) & 1.9 & 3.5 & 4.8 & 7.1\\
\bottomrule
\end{tabular}
\end{center}

À dix fois son capital, le book ne se trouve jamais à court de collatéral. Il encaisse le mouvement et en récupère presque tout. À cinquante fois, le book se situe déjà à
son plafond de collatéral, égal à un sur 2~\%, si bien que même les journées ordinaires imposent de petites ventes~: 25 jours avant le stress. Pendant le stress, la décote et la marge plus élevées plafonnent son
effet de levier à 16.7, et son nominal passe de 60.7 à 17.5, vendu dans une base sous tension. Sur l'ensemble du stress et de la reprise, son P\&L rapporté au capital s'établit à $-6.6~\%$ au titre des valorisations, $-3.3~\%$ au titre des coûts des ventes forcées, $-0.6~\%$ au titre du
repo plus cher, et $+0.8~\%$ au titre du spread encore perçu. Les valorisations se sont redressées pour les positions conservées, et non pour celles qu'il a vendues
(le \cref{fig:s2:the-bond-futures-basis-trade:capital}).

\begin{omfigure}
\begin{tikzpicture}
\begin{axis}[width=11cm,height=5.4cm,xlabel={\small année},ylabel={\small capital (départ 1)},tick label style={font=\footnotesize},
  xmin=0,xmax=5,ymin=0.95,ymax=1.45,grid=major,grid style={gray!25},table/col sep=comma,
  legend style={font=\scriptsize,at={(0.03,0.97)},anchor=north west}]
\addplot[omDef,thick] table[x=year,y=L10]{figdata/strategies-2/11-the-bond-futures-basis-trade/capital.csv};
\addplot[gray,thick] table[x=year,y=L20]{figdata/strategies-2/11-the-bond-futures-basis-trade/capital.csv};
\addplot[omThm,thick] table[x=year,y=L50]{figdata/strategies-2/11-the-bond-futures-basis-trade/capital.csv};
\legend{10$\times$,20$\times$,50$\times$}
\end{axis}
\end{tikzpicture}

\omcaption{Capital du book de basis synthétique à trois effets de levier sur cinq ans, avec une tension de financement à l'année 2.5. Les pertes du livre à cinquante fois
sont figées par des ventes forcées~; le livre à dix fois reprend le basis à son retour. Données~: \texttt{s2\_basis.paths}.}\label{fig:s2:the-bond-futures-basis-trade:capital}
\end{omfigure}

\section{Mesurer l'encombrement}

Le risque de la transaction ne réside dans aucun livre en particulier~: il est la somme de ceux de tous. Les mesures publiques comprennent les positions vendeuses sur contrats à terme des fonds à effet de levier dans le rapport de la CFTC, les prêts de repo des dealers,
et l'écart entre le rendement implicite des contrats à terme et celui du comptant. Aucune ne distingue les basis trades des autres positions. Barth et Kahn ont utilisé
des déclarations réglementaires qui ne sont pas publiques. Un trader qui gère le livre doit mesurer ce que le marché devrait absorber si le collatéral de chaque détenteur
augmentait simultanément. Le choc synthétique montre à quelle vitesse ce plafond de collatéral devient contraignant lorsque cela se produit.

\section{Fiches de stratégie}

\begin{strategyfile}{Base cash-futures longue}\label{strat:s2:the-bond-futures-basis-trade:long}
\sfield{Qui vous paie, et pourquoi} Les sociétés de gestion qui veulent de la duration par les contrats à terme sans mobiliser leur bilan, et qui poussent les contrats à terme
au-dessus de la fair value.
\sfield{Instruments et plateformes} Obligations du Trésor achetées et financées en repo~; contrats à terme sur obligations du Trésor vendus au CBOT.
\sfield{Signal} Taux de repo implicite comparé au taux de repo du moins cher à livrer (CTD).
\sfield{Dimensionnement et exécution} Effet de levier nettement inférieur à un sur (décote plus marge) dans des conditions de tension~; échéance du repo alignée sur celle de la transaction.
\sfield{Coûts} Spread du repo, frais sur contrats à terme, coûts de roll.
\sfield{Comment elle meurt} Une hausse de la décote et de la marge qui force des ventes dans une base mouvante~; un repo qui ne peut pas être renouvelé.
\sfield{Horizon, capacité, infrastructure} Rolls trimestriels~; lignes de repo et de compensation des contrats à terme.
\sfield{Backtest honnête} Tester le collatéral sous contrainte, pas seulement les prix.
\sfield{Sources} Barth et Kahn (2021)~: 400 à 500 milliards de dollars au plus haut~; ce chapitre~: 10.3~\% par an à cinquante fois et $-9.6~\%$ après une tension de financement.
\end{strategyfile}

\begin{strategyfile}{Trade d'option de livraison}\label{strat:s2:the-bond-futures-basis-trade:delivery}
\sfield{Qui vous paie, et pourquoi} Les acheteurs de contrats à terme qui laissent au vendeur le choix de l'obligation et du moment.
\sfield{Instruments et plateformes} Obligations livrables~; le future.
\sfield{Signal} Le net basis par rapport à la valeur de l'option d'échange à mesure que les taux et la courbe se déplacent.
\sfield{Dimensionnement et exécution} Acheter le net basis des obligations susceptibles de devenir les moins chères à livrer.
\sfield{Coûts} Fourchette sur les obligations off-the-run.
\sfield{Comment elle meurt} Les taux restent à un niveau tel que l'obligation la moins chère ne change pas.
\sfield{Horizon, capacité, infrastructure} Un cycle de livraison~; un modèle d'option de livraison.
\sfield{Backtest honnête} Paniers de livraison effectifs et facteurs de concordance.
\sfield{Sources} Aucune performance vérifiée.
\end{strategyfile}

\begin{strategyfile}{Signal d'encombrement du basis trade}\label{strat:s2:the-bond-futures-basis-trade:crowding}
\sfield{Qui vous paie, et pourquoi} Les vendeurs forcés lors d'un débouclage~: un trader disposant de capital libre achète le basis qu'ils vendent.
\sfield{Instruments et plateformes} Comme pour la base longue.
\sfield{Signal} Positions vendeuses sur contrats à terme des fonds à effet de levier à des extrêmes, conditions du collatéral en hausse, basis qui s'élargit rapidement.
\sfield{Dimensionnement et exécution} Position réduite, ajoutée au moment où les ventes forcées culminent~; financée par du repo à terme.
\sfield{Coûts} Marchés larges en période de tension.
\sfield{Comment elle meurt} Achat trop tôt dans un débouclage plus long.
\sfield{Horizon, capacité, infrastructure} Quelques semaines~; financement engagé.
\sfield{Backtest honnête} Peu d'épisodes~; la catégorie de la CFTC mélange les stratégies.
\sfield{Sources} CFTC Traders in Financial Futures (série dérivée de ce chapitre)~; aucune performance vérifiée.
\end{strategyfile}

\section{Tutoriel~: cinquante fois}\label{tut:s2:the-bond-futures-basis-trade:tut}

\begin{tutorial}
\textbf{Objectif.} Faire tourner le book de basis à plusieurs effets de levier à travers une tension de financement, et lire les positions des fonds à effet de levier dans les données de la CFTC.
\textbf{État final~:} le tableau et les deux figures.
\begin{tutsteps}
\item \textbf{Le livre}.
\omcode{code/firm/basistrade/firm_basistrade.py}{66}{96}{Portage, valorisations, financement et ventes forcées, avec réendettement lors des rolls.}\label{lst:s2:the-bond-futures-basis-trade:book}
\item \textbf{Par effet de levier}.
\omcode{code/strategies-2/11-the-bond-futures-basis-trade/python/s2_basis.py}{37}{50}{Rendement avant, creux et perte après la tension, par effet de levier.}\label{lst:s2:the-bond-futures-basis-trade:table}
\item \textbf{Exécuter} \texttt{s2\_fetch\_tff.py} une fois, puis \texttt{table()}, \texttt{parts()} et \texttt{fig\_basis.py}.
\end{tutsteps}
\textbf{Pour aller plus loin.} Financer avec du repo à terme qui ne peut pas être retiré pendant la tension, et déterminer l'effet de levier qu'il autorise~; conserver un coussin de liquidités
de 2~\% et comparer~; faire survenir la tension deux fois dans l'année.
\end{tutorial}

\section{Construction~: basis trade}\label{bld:s2:the-bond-futures-basis-trade:basistrade}

\begin{build}
\bfield{Rôle} Un book de basis cash-futures à effet de levier, avec décotes, marge sur contrats à terme, coût du repo et désendettement forcé.
\bfield{Interface} \texttt{BasisConfig(\dots)}, \texttt{simulate\_basis(cfg)}, \texttt{run\_book(sim, cfg, leverage)}.
\bfield{Règles} Le capital couvre la décote plus la marge sur le nominal détenu, sinon le livre vend~; pas de réendettement pendant la tension~; réendettement trimestriel ensuite.
\bfield{Tests d'acceptation} \texttt{code/firm/basistrade/tests/}~: un livre calme gagne le spread multiplié par l'effet de levier~; ventes forcées seulement lorsque le capital
tombe sous le seuil, et le capital couvre ensuite le collatéral~; la forme du scénario de tension.
\bfield{Pour aller plus loin} Repo à terme~; l'option de livraison~; plusieurs contrats.
\end{build}

\begin{omsources}
\item D.~Barth et R.~J.~Kahn, ``Hedge funds and the Treasury cash-futures disconnect'', Office of Financial Research Working Paper 21-01, 2021.
\item Board of Governors of the Federal Reserve System, communiqués de presse des 15 et 23 mars 2020.
\item CFTC, Traders in Financial Futures, 2010--2026.
\end{omsources}

\section{Exercices}

\begin{exercise}[$\star$]\label{exo:s2:the-bond-futures-basis-trade:1}
La décote et la marge sur contrats à terme valent chacune 1~\%. Quel est le levier maximal~? Et à 3~\% chacune~?
\end{exercise}

\begin{exercise}[$\star$]\label{exo:s2:the-bond-futures-basis-trade:2}
Un livre rapporte 20 points de base par an sur son nominal, à cinquante fois son capital. Que rapporte-t-il sur son capital, avant coûts~?
\end{exercise}

\begin{exercise}[$\star$]\label{exo:s2:the-bond-futures-basis-trade:3}
Le basis évolue de 60 centimes pour 100 contre un livre à cinquante fois. Que perd-il sur son capital, avant toute vente~?
\end{exercise}

\begin{exercise}[$\star\star$]\label{exo:s2:the-bond-futures-basis-trade:4}
Pourquoi le livre à dix fois se redresse-t-il alors que le livre à cinquante fois ne se redresse pas~?
\end{exercise}

\begin{exercise}[$\star\star$]\label{exo:s2:the-bond-futures-basis-trade:5}
De combien en pourcentage les hedge funds ont-ils réduit leurs positions vendeuses sur contrats à terme à 2, 5 et 10 ans entre le 18 février et le 17 mars 2020~?
\end{exercise}

\begin{exercise}[$\star\star$]\label{exo:s2:the-bond-futures-basis-trade:6}
Pourquoi la série des fonds à effet de levier de la CFTC n'est-elle qu'une mesure approximative du basis trade~?
\end{exercise}

\begin{exercise}[$\star\star\star$]\label{exo:s2:the-bond-futures-basis-trade:7}
\emph{Programmation.} Relancer le livre à cinquante fois avec \texttt{BasisConfig(haircut\_stress=0.01, margin\_stress=0.01)}, de sorte que les conditions du collatéral
ne changent pas. Que se passe-t-il, et que cela montre-t-il~?
\end{exercise}

\begin{exercise}[$\star\star\star$]\label{exo:s2:the-bond-futures-basis-trade:8}
\emph{Trouver la faille.} «~La base converge toujours à la livraison, donc la transaction ne peut pas perdre si l'on tient la position jusqu'à la livraison.~»
\end{exercise}

\section{Problème~: Fifty Times}

\begin{problem}[{Problème du week-end --- un livre de basis trade à levier élevé}]\label{pb:s2:the-bond-futures-basis-trade:1}
Le livre synthétique du chapitre, les données de la CFTC et le registre public.

\medskip\noindent\textbf{Partie I --- Mécanique.}
\begin{enumerate}
\item Définir la base brute, la base nette et le taux de repo implicite.
\item Quand une position longue sur la base rapporte-t-elle~?
\item Définir l'\omterm{def:s2:the-bond-futures-basis-trade:leverage}{effet de levier du basis trade} et le \omterm{def:s2:the-bond-futures-basis-trade:rollover}{risque de renouvellement du repo}.
\item Qu'est-ce qui limite l'effet de levier~?
\end{enumerate}

\medskip\noindent\textbf{Partie II --- Qui le pratique.}
\begin{enumerate}[resume]
\item Pourquoi les sociétés de gestion poussent-elles les contrats à terme au-dessus de la fair value~?
\item Qu'ont estimé Barth et Kahn au sujet de la taille du marché~?
\item Que montre la série de la CFTC~?
\item Pourquoi les auteurs qualifient-ils les hedge funds d'entrepôts~?
\end{enumerate}

\medskip\noindent\textbf{Partie III --- Mars 2020.}
\begin{enumerate}[resume]
\item Que s'est-il passé sur les marchés du Trésor et du repo~?
\item Que firent les hedge funds~?
\item Qu'annonça la Réserve fédérale~?
\item Les basis trades furent-ils la cause principale, selon Barth et Kahn~?
\end{enumerate}

\medskip\noindent\textbf{Partie IV --- Le verdict.}
\begin{enumerate}[resume]
\item Énoncer le \emph{résultat nommé}~: le rendement du capital de la transaction selon l'effet de levier et la perte dans le scénario de tension de financement prévu.
\item Décomposer la perte du livre à cinquante fois.
\item Que changerait le financement à terme~?
\item Comment mesureriez-vous l'encombrement~?
\item Comment dimensionneriez-vous le livre~?
\item Comment feriez-vous un backtest honnête de celui-ci~?
\item Quelle stratégie du fichier tire profit des ventes forcées des autres~?
\item En une phrase~: que apporte un basis trader au marché~?
\end{enumerate}
\end{problem}

\section{Questions d'entretien}

\begin{interviewq}[$\star$ \iqroles{trader}]\label{iq:s2:the-bond-futures-basis-trade:1}
Expliquez le basis trade entre le comptant et les contrats à terme sur le Trésor.
\end{interviewq}

\begin{interviewq}[$\star\star$ \iqroles{risque}]\label{iq:s2:the-bond-futures-basis-trade:2}
Un livre de basis est à un levier de 50 fois. Quel scénario de tension feriez-vous tourner, et quelle limite fixeriez-vous~?
\end{interviewq}

\begin{interviewq}[$\star\star$ \iqroles{chercheur}]\label{iq:s2:the-bond-futures-basis-trade:3}
Comment calculeriez-vous le taux de repo implicite, et que vous apprend son écart au repo~?
\end{interviewq}

\begin{interviewq}[$\star\star$ \iqroles{trader}]\label{iq:s2:the-bond-futures-basis-trade:4}
Votre prêteur en repo relève la décote de 1~\% à 3~\% pendant la nuit. Que faites-vous~?
\end{interviewq}

\begin{interviewq}[$\star\star$ \iqroles{développeur}]\label{iq:s2:the-bond-futures-basis-trade:5}
Concevez le processus quotidien d'un livre de basis~: données, collatéral, appels de marge, roll.
\end{interviewq}

\begin{interviewq}[$\star\star\star$ \iqroles{chercheur}]\label{iq:s2:the-bond-futures-basis-trade:6}
Montrez que, avec un capital $C$, un nominal $N$ et un taux de rémunération du collatéral $k$ par unité de nominal, une perte de $x$ par unité de nominal force des ventes lorsque $C - Nx < Nk$, et déterminez
le mouvement le plus grand qu'un livre à l'effet de levier $L$ peut subir sans vendre.
\end{interviewq}
